% Source: kactl

\chapter{Geometry}

\code{Geometry Primitives}{geometry/primitives.cpp}

\codeWithDescription{Point 2D}{geometry/Point.h}
[Class to handle points in the plane. T can be e.g. double or long long. (Avoid int.)]

\codeWithDescription{Line Intersection}{geometry/lineIntersection.h}[][]
% [If a unique intersection point of the lines going through s1,e1 and s2,e2 exists \{1, point\} is returned. 
% If no intersection point exists \{0, (0,0)\} is returned and if infinitely many exists \{-1, (0,0)\} is returned. 
% The wrong position will be returned if P is Point<ll> and the intersection point does not have integer coordinates. 
% Products of three coordinates are used in intermediate steps so watch out for overflow if using int or ll.]%
% [graphics/geometry/lineIntersection]
[
    auto res = lineInter(s1,e1,s2,e2);
 	if (res.first == 1)
		cout << "intersection point at " << res.second << endl;
]

\code{Line Equation}{geometry/line_equation.cpp}
\codeWithDescription{Segment Distance}{geometry/SegmentDistance.h}[]
% [Returns the shortest distance between point p and the line segment from point s to e.]
[]
[
    Point<double> a, b(2,2), p(1,1);
 	bool onSegment = segDist(a,b,p) < 1e-10;
]

\codeWithDescription{Segment Intersection}{geometry/SegmentIntersection.h}[][]
% [
% If a unique intersection point between the line segments going from s1 to e1 and from s2 to e2 exists then it is returned.
% If no intersection point exists an empty vector is returned.
% If infinitely many exist a vector with 2 elements is returned, containing the endpoints of the common line segment.
% The wrong position will be returned if P is Point<ll> and the intersection point does not have integer coordinates.
% Products of three coordinates are used in intermediate steps so watch out for overflow if using int or long long.
% ]
% [graphics/geometry/SegmentIntersection]
[vector<P> inter = segInter(s1,e1,s2,e2);
 if (sz(inter)==1)
   cout << "segments intersect at " << inter[0] << endl;
]

\codeWithDescription{Side Of}{geometry/sideOf.h}
% [
%  Returns where $p$ is as seen from $s$ towards $e$. 1/0/-1 $\Leftrightarrow$ left/on line/right.
%  If the optional argument $eps$ is given 0 is returned if $p$ is within distance $eps$ from the line.
%  P is supposed to be Point<T> where T is e.g. double or long long.
%  It uses products in intermediate steps so watch out for overflow if using int or long long.
% ][]
[][]
[bool left = sideOf(p1,p2,q)==1;]

\codeWithDescription{On Segment}{geometry/OnSegment.h}
[Returns true iff p lies on the line segment from s to e.]
[]
[\texttt{(segDist(s,e,p)<=epsilon)} instead when using Point<double>.]

\codeWithDescription{Linear Transformation}{geometry/linearTransformation.h}
[Apply the linear transformation (translation, rotation and scaling) which takes line p0-p1 to line q0-q1 to point r.]
% [graphics/geometry/linearTransformation]

% \codeWithDescription{Circle Intersection}{geometry/CircleIntersection.h}
% [
%   Computes the pair of points at which two circles intersect. Returns false in case of no intersection.
% ]
\code{Circle Line Intersection}{geometry/CircleLineIntersection.cpp}
\code{Circle Intersection}{geometry/CircleIntersection.h}
\code{Circle Polygon Intersection}{geometry/CirclePolygonIntersection.h}
\code{Circle Line}{geometry/CircleLine.h}

\codeWithDescription{Circle Tangents}{geometry/CircleTangents.h}
[Finds the external tangents of two circles, or internal if r2 is negated.
 Can return 0, 1, or 2 tangents -- 0 if one circle contains the other (or overlaps it, in the internal case, or if the circles are the same);
 1 if the circles are tangent to each other (in which case .first = .second and the tangent line is perpendicular to the line between the centers).
 .first and .second give the tangency points at circle 1 and 2 respectively.
 To find the tangents of a circle with a point set r2 to 0. $O(n)$
 ]

 \codeWithDescription{Minimum Enclosing Circle}{geometry/MinimumEnclosingCircle.h}
 [Computes the minimum circle that encloses a set of points. $O(n)$]

 \codeWithDescription{Inside Polygon}{geometry/InsidePolygon.h}
 [
 Returns true if p lies within the polygon. If strict is true, it returns false for 
 points on the boundary. The algorithm uses products in intermediate steps so watch out for overflow. $O(n)$
 ]
 []
 [vector<P> v = {P{4,4}, P{1,2}, P{2,1}};
 bool in = inPolygon(v, P{3, 3}, false);
 ]

\codeWithDescription{Polygon Area}{geometry/PolygonArea.h}
[Returns twice the signed area of a polygon. Clockwise enumeration gives negative area. 
Watch out for overflow if using int as T!]

\codeWithDescription{Polygon Center}{geometry/PolygonCenter.h}
[Returns the center of mass for a polygon. Time: $O(n)$]

\code{Polygon Cut}{geometry/PolygonCut.h}
\code{Polygon Union}{geometry/PolygonUnion.h}

\code{Extreme Vertex}{geometry/extremeVertex.cpp}

\code{Manhattan MST}{geometry/ManhattanMST.h}

\codeWithDescription{Convex Hull}{geometry/ConvexHull.h}
[Returns a vector of the points of the convex hull in counter-clockwise order.
Points on the edge of the hull between two other points are not considered part of the hull. $O(nlogn)$
]

\codeWithDescription{Hull Diameter}{geometry/HullDiameter.h}
[Returns the two points with max distance on a convex hull (ccw, no duplicate/collinear points). $O(n)$]

\codeWithDescription{Point Inside Hull}{geometry/PointInsideHull.h}
[Determine whether a point t lies inside a convex hull (CCW
 order, with no collinear points). Returns true if point lies within
 the hull. If strict is true, points on the boundary aren't included. $O(logN)$]

\codeWithDescription{Line Hull Intersection}{geometry/LineHullIntersection.h}
[Line-convex polygon intersection. The polygon must be ccw and have no collinear points.
 lineHull(line, poly) returns a pair describing the intersection of a line with the polygon:
  \begin{itemize}
    \item $(-1, -1)$ if no collision,
    \item $(i, -1)$ if touching the corner $i$,
    \item $(i, i)$ if along side $(i, i+1)$,
    \item $(i, j)$ if crossing sides $(i, i+1)$ and $(j, j+1)$.
  \end{itemize}
  In the last case, if a corner $i$ is crossed, this is treated as happening on side $(i, i+1)$.
  The points are returned in the same order as the line hits the polygon.
 \texttt{extrVertex} returns the point of a hull with the max projection onto a line. $O(logN)$]

\code{Half Plane Intersection}{geometry/halfPlaneIntersection.cpp}
\code{Minkowski Sum & Dif}{geometry/minkowski.cpp}
\code{Maximum Inscribed Circle}{geometry/maxInscribedCircle.cpp}
\code{Rotating Calipers}{geometry/rotatingCalipers.cpp}

\codeWithDescription{Closest Pair}{geometry/ClosestPair.h}
[Finds the closest pair of points. O(nlogn)]